% =========================================================
% Proof Stub: Regrouping and Convergence
% Source: volume-iii/analysis/sequences/notes/series/notes-series-rearrangements.tex
% =========================================================

\clearpage
\phantomsection
\hypertarget{proof-RA-ABB-T-regrouping}{}

\subsubsection[Regrouping and Convergence]{Proof --- Regrouping and Convergence}

\bigskip

\noindent
\textbf{Source.} See \texttt{notes-series-rearrangements.tex}.

\vspace{0.75em}

\noindent
\textbf{Goal.} If $\sum a_n = S$ converges, then any regrouping (without reordering) also converges to $S$.

\vspace{0.75em}

\noindent
\textbf{Logical form.}
\[
  \sum a_n = S \Rightarrow \sum_{k=1}^\infty \left(\sum_{n=n_k}^{n_{k+1}-1} a_n\right) = S.
\]

\vspace{0.75em}

\noindent
\textbf{Proof strategy.} The partial sums of the regrouped series are a subsequence of the original partial sums $(s_n)$. Subsequences of a convergent sequence converge to the same limit.

\vspace{0.75em}

\noindent
\textbf{Proof.}
\begin{proof}
\Claim If $\sum a_n = S$ converges, then any regrouping (without reordering) also converges to $S$.

\Given 

\Goal 

\Strategy The partial sums of the regrouped series are a subsequence of the original partial sums $(s_n)$. Subsequences of a convergent sequence converge to the same limit.

\medskip

% --- grouped partial sums are s_{n_1}, s_{n_2}, s_{n_3}, ... (a subsequence of (s_n)) ---
% --- s_n → S ⇒ every subsequence → S ---
% --- so grouped series converges to S ---

\end{proof}

\begin{remark*}[Proof shape]
Grouping selects a subsequence of partial sums; subsequences of convergent sequences inherit the limit.
\end{remark*}

\begin{remark*}[Dependencies]
\begin{itemize}
  \item Subsequences Inherit Limits.
  \item The regrouped partial sums are a subsequence of $(s_n)$.
\end{itemize}
\end{remark*}
